% ==============================================================================
% 1. FONTS & BEAMER MATH CONFIGURATION
% ==============================================================================
\usefonttheme{professionalfonts}
\setbeamerfont{table}{family=\sffamily}
\usepackage[defaultmathsizes]{mathastext}
\usepackage{microtype}

% ==============================================================================
% 1. FONTS & BEAMER MATH CONFIGURATION
% ==============================================================================
\usefonttheme{professionalfonts}
\renewcommand{\familydefault}{\sfdefault}
\usepackage[defaultmathsizes]{mathastext}
\usepackage{microtype}

% ==============================================================================
% 2. UTILITIES & CONVENIENCE MACROS
% ==============================================================================
\usepackage{csquotes}% load before biblatex
\usepackage{xspace}
\usepackage{standalone}
\usepackage{appendixnumberbeamer}% do not count appendix frames in total frame count

% ==============================================================================
% 3. UNITS & CHEMISTRY
% ==============================================================================
\usepackage{siunitx}
\sisetup{
    mode=text,
    reset-math-version = false,
    range-phrase={--},
    range-units=single,
    per-mode=symbol,
}
\usepackage[version=4]{mhchem}

% ==============================================================================
% 4. GRAPHICS, TIKZ & PGFPLOTS
% ==============================================================================
\usepackage{caption}
\usepackage{tikz}
\usetikzlibrary{shadows,spy,calc}
\usepackage{pgfplots}
\pgfplotsset{compat=1.18}
\usepackage{pgfplotstable}
\pgfplotstableset{
    col sep=comma,
    every column/.style={string type},
    every head row/.style={before row={\toprule}},
    every row no 0/.style={after row=\midrule},
    every last row/.style={after row=\bottomrule}
}
\usepackage{shadowtext}
\shadowoffset{1pt}
\shadowcolor{ubRed}
\usepackage{animate}
\usepackage{qrcode}
\usepackage{ccicons}
\usepackage{fontawesome5}

% ==============================================================================
% 5. TABLES & CODE LISTINGS
% ==============================================================================
\usepackage{booktabs}
\usepackage{colortbl}% load before listings to prevent macro conflicts
\usepackage{multirow}
\usepackage{listings}
\lstset{basicstyle=\scriptsize\ttfamily,breaklines=true}
\usepackage{lstlinebgrd}
% Compatibility fix for lstlinebgrd with LaTeX kernels >= 2021-11-15
\makeatletter
\let\old@lst@linebgrd\lst@linebgrd
\def\lst@linebgrd{\lst@basicstyle\old@lst@linebgrd}
\makeatother

% ==============================================================================
% 6. BIBLIOGRAPHY & CITATIONS
% ==============================================================================
\usepackage[%
    backend=biber,%
    sorting=none,%
    backref=true,%
    backrefstyle=none,%
    citetracker=true,%
    isbn=false,%
    url=false,%
    maxnames=1,minnames=1,% Show only first author in references
]{biblatex}
% show "cited on slide(s) <numbers>" in bibliography
\DefineBibliographyStrings{english}{%
    backrefpage={slide},%
    backrefpages={slide}%
}
\renewbibmacro*{pageref}{%
    \iflistundef{pageref}%
    {}%
    {\printtext[parens]{\bibstring{backrefpages}\addspace\printlist{pageref}}}%
}
\renewbibmacro{in:}{}% no "In" in bibliography
\DeclareSourcemap{%
    \maps[datatype=bibtex]{%
        \map{%
            \step[fieldset=journal,null]%
            \step[fieldset=editor,null]%
            \step[fieldset=publisher,null]%
            \step[fieldset=series,null]%
            \step[fieldset=volume,null]%
            \step[fieldset=number,null]%
            \step[fieldset=pages,null]%
        }%
    }%
}
\addbibresource{../../../../Documents/library.bib}

% ==============================================================================
% 7. VERSION CONTROL & METADATA
% ==============================================================================
\usepackage[local]{gitinfo2}
\usepackage[%
    english,%
    showseconds=false,%
    showzone=false%
]{datetime2}
\DTMsettimestyle{iso}
\usepackage{appendixnumberbeamer}% do not count appendix frames in total frame count
\usepackage{etoolbox}

% ==============================================================================
% 8. CUSTOM MACROS & COLOR DEFINITIONS
% ==============================================================================
\newlength\imagescale% needed for scalebars
\newcommand{\uct}{{\textmu}CT\xspace}
\newcommand{\eg}{e.\,g.\xspace}
\newcommand{\ie}{i.\,e.\xspace}
% Image dimensions — use \renewcommand{\imagewidth}{...} on slides that need a different size
\newcommand{\imagewidth}{\linewidth}
\newcommand{\standardimagewidth}{\linewidth}
\newcommand{\imageheight}{0.72\paperheight}
\newcommand{\standardimageheight}{0.72\paperheight}

% Convenience functions to "cite" a source image.
% Place the commands immediately after \includegraphics, flush right
\newlength{\sourcegap}
\setlength{\sourcegap}{0.25ex}
\newcommand{\source}[2]{% \source{example.com/image}{CC BY 4.0}
    % Cite a (short) link directly
    \raisebox{-\height-\sourcegap}{%
        \makebox[0pt][r]{%
            \tiny
            \href{https://#1}{#1}%
            \ifblank{#2}{}{ #2}%
        }%
    }%
}
\newcommand{\sourcecite}[2]{% \sourcecite{BibTeX key}{Fig. 3}
    % Cite a reference
    \raisebox{-\height-\sourcegap}{\makebox[0pt][r]{%
        \tiny
        From~\cite{#1}
        \ifblank{#2}{}{, #2}%
        }%
    }%
}
\newcommand{\sourcelink}[3]{% \sourcelink{https://example.com/very/long/url}{example.com}{CC BY-SA}
    % Cite a link
    \raisebox{-\height-\sourcegap}{%
        \makebox[0pt][r]{%
            \tiny
            \href{#1}{#2}%
            \ifblank{#3}{}{ #3}%
        }%
    }%
}
\newcommand{\sourceqr}[1]{% \sourceqr{https://example.com/some/very/long/source/url}
    % Cite (an image from) a (long) URL with a QR code
    \makebox[0pt][l]{\hspace{\sourcegap}\raisebox{\depth}{\qrcode[height=1cm]{#1}}}%
}
\newcommand{\citeslide}[5]{% slideimage, lecturelink, lecturetitle, author, page
    \fboxsep=0pt%
    \fcolorbox{ubRed!62}{white}{\includegraphics[%
        width=\linewidth,height=\imageheight,keepaspectratio]{#1}}%
    \raisebox{-\height}{\makebox[0pt][r]{%
        \tiny\href{#2}{\emph{#3} by #4}, Slide #5}}%
}

% Complementary to ubRed (#e4003c), defined in the theme .sty
\definecolor{ubRedComplementary}{HTML}{00E3A6}

% Notes always on white, regardless of the notes mode
\setbeamercolor{note page}{bg=white}
\setbeamercolor{note date}{bg=white}

\tikzset{shadowed/.style={preaction={transform canvas={shift={(1pt,-1pt)}},draw=ubRed}}}
\tikzset{every picture/.style={thick}}% globally thicker lines

% ==============================================================================
% 9. BEAMER PROGRESS BAR & NAVIGATION TEMPLATES
% ==============================================================================
% Progress bar marker
\usepackage{xfp} % For robust math calculations
% Setup dimensions
\newlength{\barstart}
\setlength{\barstart}{15.2mm} % beamerouterthemeubbeamer2023.sty:\setbeamersize{text margin left=15.2mm}
\newlength{\barend}
\setlength{\barend}{\paperwidth - 6mm} % beamerouterthemeubbeamer2023.sty: \put(15.2,-1.36){\usebeamercolor[fg]{ub page rule}\rule{\paperwidth-21.2mm}{1.36mm}}
\newlength{\barmaxwidth}
\setlength{\barmaxwidth}{\barend - \barstart}

% Total frame count for the bar.
% \appendix is patched to record the frame number into the .aux
% Read back on the next run; until then the bar simply isn't drawn.
\providecommand{\appendixframeno}{1}
\makeatletter
\AtBeginDocument{%
    \let\pb@origappendix\appendix
    \renewcommand{\appendix}{%
        \immediate\write\@auxout{\gdef\string\appendixframeno{\the\c@framenumber}}%
    \pb@origappendix}%
}
\makeatother

% Break marker: \saveframenumber writes the frame number to \jobname.mkr; on the
% next run it is read back and drawn as a white tick on the bar.
% Use \saveframenumber *after* the slide where you'd like to make a break.
% This is marked with a white line in the top bar.
\makeatletter
\newwrite\@markerout
\providecommand{\myframenumber}{-1}% -1 = no marker
\AtBeginDocument{\InputIfFileExists{\jobname.mkr}{}{}}
\newcommand{\saveframenumber}{%
    \immediate\openout\@markerout=\jobname.mkr%
    \immediate\write\@markerout{\gdef\string\myframenumber{\insertframenumber}}%
    \immediate\closeout\@markerout%
}
\makeatother

% The actual progress bar command
\makeatletter
\newcommand{\myprogressbar}{%
    \ifnum\insertframenumber>2
    % Prefer the frozen pre-appendix count; fall back to beamer's own total if there
    % is no \appendix. Either way, don't draw until a plausible total is known.
    \edef\pb@total{\appendixframeno}%
    \ifnum\pb@total>3\else\edef\pb@total{\inserttotalframenumber}\fi
    \ifnum\pb@total>3
    \begin{tikzpicture}[remember picture, overlay]
        % Calculate Ratio using xfp (Current - 3) / (End - 3)
        % We want to start the marker at frame three, after UB logo and title frame
        \edef\ratio{\fpeval{min(max((\insertframenumber - 3) / max(\pb@total - 3, 1), 0), 1)}}
        % Break marker first, so the position circle sits on top of it
        \ifnum\myframenumber>2
        \edef\markratio{\fpeval{min(max((\myframenumber - 3) / max(\pb@total - 3, 1), 0), 1)}}
        \draw[white, line width=0.4mm] (\barstart + \markratio\barmaxwidth, 0) -- ++(0, -1.36mm);
        \fi
        % Draw a circle as marker in the red top bar
        \fill[white] (\barstart + \ratio\barmaxwidth, -0.68mm) circle (0.68mm); % 0.68 is 1.36/2
    \end{tikzpicture}
    \fi
    \fi
}
\makeatother

% Set the progress bar as header, only in presentation mode
\mode<beamer>{\setbeamertemplate{headline}{\myprogressbar}}

% Show section ToC at the start of each section (presentation mode only)
\mode<beamer>{%
    \AtBeginSection[]{%
        \begin{frame}{Contents}%
            \tableofcontents[currentsection,currentsubsection,hideothersubsections]%
        \end{frame}%
    }%
}%
